\documentclass[10pt]{article}
\usepackage[spanish,es-noshorthands,es-nolayout]{babel}
\newcommand{\nombre}{Luis Valverde Ramos}
\newcommand{\pie}{Currículum, noviembre de 2026}
% Diseño compartido de cv-es.tex y cv-en.tex. Compila con pdfLaTeX (también
% en Overleaf): Archivo es la misma letra de la web y el verde es su acento.
% Las referencias van en privado/referencias-<idioma>.tex, fuera del repo, y
% solo entran si se compila con \def\conreferencias{} (build.py --privado).

\usepackage[a4paper,left=1.8cm,right=1.8cm,top=1.4cm,bottom=1.6cm,footskip=0.8cm]{geometry}
\usepackage[T1]{fontenc}
\usepackage[lining,tabular,mainfont,semibold]{Archivo}
\usepackage{microtype}
\usepackage{xcolor}
\usepackage{enumitem}
\usepackage{needspace}
\usepackage{fancyhdr}
\usepackage{lastpage}
\usepackage[hidelinks]{hyperref}

\definecolor{acento}{HTML}{2E6B4F}
\definecolor{tinta}{HTML}{191E1B}
\definecolor{gris}{HTML}{525A54}
\definecolor{linea}{HTML}{C2C9BF}

\setlength{\parindent}{0pt}
\linespread{1.05}
\frenchspacing
\brokenpenalty=10000 % ninguna palabra partida entre dos páginas
\AtBeginDocument{\color{tinta}}

% Columna de fechas a la izquierda y texto a la derecha
\newlength{\fechas}\setlength{\fechas}{2.45cm}
\newlength{\separa}\setlength{\separa}{0.4cm}

% Pie: nombre, fecha de actualización y página
\pagestyle{fancy}
\fancyhf{}
\renewcommand{\headrulewidth}{0pt}
\fancyfoot[C]{\scriptsize\color{gris}\nombre\ · \pie\ · \thepage/\pageref*{LastPage}}

% Cabecera: nombre, una línea de presentación y los contactos
\newcommand{\cabecera}[2]{%
  {\fontsize{23}{27}\selectfont\bfseries\nombre}\par\vspace{4pt}
  {\color{gris}#1}\par\vspace{3pt}
  {\small #2}\par\vspace{7pt}
  {\color{acento}\hrule height 1.2pt}\vspace{2pt}}
\newcommand{\enlace}[2]{\href{#1}{\color{acento}#2}}

% Sección: título en versalitas verdes con una raya fina; nunca sola al pie
\newcommand{\seccion}[1]{%
  \par\Needspace{5\baselineskip}\vspace{9pt}%
  {\color{acento}\bfseries\footnotesize\textls[110]{\MakeUppercase{#1}}}\par
  \vspace{2pt}{\color{linea}\hrule height 0.5pt}\vspace{6pt}}

% Lista de entradas con la fecha en la columna izquierda
% \leavevmode antes del color, para que la fecha comparta la línea base del
% cargo; \smash, para que una fecha de dos líneas no empuje el texto
\newcommand{\fecha}[1]{\smash{\parbox[t]{\fechas}{\raggedright\small\leavevmode\color{gris}#1}}}
\newlist{entradas}{description}{1}
\setlist[entradas]{leftmargin=\dimexpr\fechas+\separa\relax, labelwidth=\fechas,
  labelsep=\separa, align=left, font=\normalfont, itemsep=5pt, topsep=0pt,
  parsep=1pt, partopsep=0pt}
% Cada entrada empieza en una página con sitio para su título y dos líneas más
\newcommand{\entrada}[2][]{\Needspace{3\baselineskip}\item[\fecha{#1}]#2}

% Cargo e institución, y un detalle opcional en gris debajo
\newcommand{\cargo}[2]{{\bfseries #1}, #2\par}
\newcommand{\detalle}[1]{{\small\color{gris}#1}\par}

% Viñetas dentro de una entrada
\newlist{puntos}{itemize}{1}
\setlist[puntos]{label={\color{acento}\textbullet}, leftmargin=1.1em,
  labelsep=0.45em, itemsep=1.5pt, topsep=2pt, parsep=0pt, before=\small}


\begin{document}

\cabecera{Politólogo (PUCP) · Lima, Perú}{%
  \enlace{mailto:luis.valverde@pucp.edu.pe}{luis.valverde@pucp.edu.pe} \ · \
  \enlace{https://luiseduardovalverderamos.github.io/es/}{luiseduardovalverderamos.github.io} \ · \
  \enlace{https://www.linkedin.com/in/luisvalverderamos/}{linkedin.com/in/luisvalverderamos}}

\seccion{Perfil}
Politólogo por la Pontificia Universidad Católica del Perú (PUCP), primer puesto de la promoción de egreso, con tesis de licenciatura aprobada con mención sobresaliente. Investigo comportamiento electoral, comportamiento legislativo e implementación de políticas públicas en el Perú con métodos mixtos y R. Soy investigador de la Escuela Electoral y de Gobernabilidad del Jurado Nacional de Elecciones, consultor del Observatorio de Corrupción del Colegio de Abogados de Lima y jefe de práctica de Economía Política en la PUCP.

\seccion{Formación}
\begin{entradas}
\entrada[2019 -- 2024]{\cargo{Licenciatura en Ciencia Política y Gobierno}{Pontificia Universidad Católica del Perú (PUCP)}
  \detalle{Primer puesto de la promoción de egreso. Tesis aprobada con mención sobresaliente (2025).}}
\entrada[2024]{\cargo{Intercambio académico en Ciencias Políticas y Sociales}{Ghent University, Bélgica}
  \detalle{Cursos de maestría: Quantitative Methods for Social Research, Public Management, Innovation Management y European Policy.}}
\entrada[2016 -- 2018]{\cargo{Diploma del Bachillerato Internacional (IB)}{Colegio de Alto Rendimiento (COAR) Amazonas}
  \detalle{Egresado destacado.}}
\end{entradas}

\seccion{Experiencia en investigación}
\begin{entradas}
\entrada[Jul. 2026 --]{\cargo{Investigador}{Escuela Electoral y de Gobernabilidad (ESEG), Jurado Nacional de Elecciones}
  \begin{puntos}
  \item Contribuyo al diseño de la agenda de investigación institucional con métodos cuantitativos y cualitativos.
  \item Trabajo en el censo de la enseñanza electoral universitaria (en curso).
  \item Enlace técnico con autoridades académicas y jurídicas en el campo de los estudios electorales.
  \end{puntos}}
\entrada[Ene. -- mar. 2026]{\cargo{Miembro voluntario}{Hechos y No Palabras, Grupo de Investigación en Estado y Sociedad (GIES), PUCP}
  \begin{puntos}
  \item Seguimiento legislativo de cuatro reformas políticas e institucionales de cara a las Elecciones Generales 2026: sistematicé votaciones nominales y construí el historial de votación de candidatos para el portal \enlace{https://hechosynopalabras.pe}{hechosynopalabras.pe}, que tuvo cerca de 80\,000 visitas.
  \end{puntos}}
\entrada[Oct. 2025 -- \newline Feb. 2026]{\cargo{Asistente de investigación}{Centro de Investigaciones Sociológicas, Económicas, Políticas y Antropológicas (CISEPA), PUCP}
  \begin{puntos}
  \item Apoyé el análisis del régimen político peruano (avance autoritario del Legislativo e inestabilidad ministerial, 2021--2025) para un artículo enviado a la \textit{Revista de Ciencia Política} (Chile).
  \item Sistematicé las votaciones del Congreso 2021--2025 para identificar coaliciones entre bancadas en temas de corte autoritario.
  \item Apoyo econométrico en R (series de tiempo interrumpidas, efectos fijos de dos vías e instrumentos \textit{shift-share}) a dos manuscritos sobre los efectos de la prohibición de la reelección subnacional en la corrupción en contrataciones públicas y en la gestión fiscal municipal, enviados a revistas internacionales indexadas.
  \item Trabajé en la construcción de un corpus de 747\,836 titulares de \textit{El Comercio} y \textit{La República} (2016--2024) para medir la cobertura mediática de cinco economías ilícitas.
  \end{puntos}}
\entrada[Sep. 2024 -- \newline Sep. 2025]{\cargo{Asistente de investigación}{Q-LAB, Laboratorio de IA y Métodos Computacionales en Ciencias Sociales, PUCP}
  \begin{puntos}
  \item En un proyecto con la Cooperación Suiza (GFP Subnacional), trabajé en un modelo de \textit{machine learning} supervisado (XGBoost y valores SHAP) sobre retrasos en la ejecución de 40\,243 obras públicas municipales (2018--2024), con registros del MEF, RENAMU e INEI; los hallazgos sirvieron de insumo para un \textit{policy brief} institucional.
  \item Construí y limpié bases de datos distritales y regionales para investigaciones sobre corrupción, desempeño municipal y migración.
  \end{puntos}}
\entrada[May. -- Ago. 2024]{\cargo{Asistente de investigación}{Daniel Encinas (Northwestern University) y Rodrigo Barrenechea}
  \begin{puntos}
  \item Procesé resultados provinciales de diez elecciones presidenciales (1980--2021) para identificar patrones geográficos del voto populista en el Perú; análisis descriptivo con datos de LAPOP.
  \end{puntos}}
\entrada[Feb. -- Dic. 2023]{\cargo{Practicante de investigación}{Instituto de Democracia y Derechos Humanos (IDEHPUCP)}
  \begin{puntos}
  \item Apoyé tres proyectos de investigación sobre gobernanza, migración y derechos humanos, en coordinación con agencias del sistema de Naciones Unidas y contrapartes del sector público.
  \item Apoyé el proceso editorial del volumen por los veinte años del Informe Final de la CVR y coescribí seis columnas para el boletín del IDEHPUCP.
  \end{puntos}}
\end{entradas}

\seccion{Consultoría}
\begin{entradas}
\entrada[Sep. 2026 --]{\cargo{Consultor}{Observatorio de Corrupción del Colegio de Abogados de Lima}
  \begin{puntos}
  \item Asesoro la primera entrega del informe del Observatorio sobre contrataciones públicas 2016--2025: evaluación de la factibilidad y la validez de indicadores construidos con registros abiertos del SEACE y el OECE.
  \end{puntos}}
\entrada[Feb. -- Ago. 2026]{\cargo{Consultor de investigación}{Liga Peruana de Debate Escolar (LPDE)}
  \begin{puntos}
  \item Investigación aplicada sobre el ecosistema de debate escolar en el Perú: diseño del estudio y análisis cuantitativo y cualitativo.
  \item Preparé el primer informe de la Liga sobre ese ecosistema (70 instituciones) y siete \textit{policy briefs} para sus decisiones estratégicas.
  \end{puntos}}
\entrada[Nov. -- Dic. 2025]{\cargo{Consultor}{ProDiálogo, por encargo del Jurado Nacional de Elecciones (JNE)}
  \begin{puntos}
  \item Integré el equipo técnico de tres instrumentos de medición del comportamiento electoral para las Elecciones Generales 2026: la Encuesta de Perfil Electoral, dos experimentos de encuesta y un mapa conductual.
  \item Participé en la segmentación del electorado peruano en cuatro perfiles, con análisis de conglomerados sobre encuestas de tres ciclos electorales (2011, 2016 y 2021), como insumo para las estrategias de educación cívica del JNE.
  \end{puntos}}
\entrada[Jul. -- Sep. 2025]{\cargo{Consultor de análisis de datos}{Sapienza Consultora, para Roche Pharma Perú}
  \begin{puntos}
  \item Análisis de registros administrativos del sector salud (2019--2025) y tres reportes para el cliente.
  \end{puntos}}
\end{entradas}

\seccion{Docencia}
\begin{entradas}
\entrada[Ago. 2026 --]{\cargo{Jefe de práctica}{Economía Política, Facultad de Ciencias Sociales, PUCP}}
\entrada[2024 -- 2025]{\cargo{Asistente docente}{Q-LAB, PUCP}
  \detalle{Introducción a R para la Investigación Social (dos ediciones) y módulo de Causal Machine Learning del curso de Econometría Aplicada Avanzada.}}
\end{entradas}

\seccion{Dirección y entrenamiento de debate}
\begin{entradas}
\entrada[2026 --]{\cargo{Director general}{Selección Nacional de Debate Escolar del Perú}
  \detalle{Dirijo y gestiono la selección: conduzco la convocatoria de debatientes y coordino con los encargados de las selecciones junior y senior.}}
\entrada[2024 --]{\cargo{Entrenador}{equipos escolares y universitarios}
  \begin{puntos}
  \item Selección Nacional de Debate Escolar (2025--2026): integré el equipo de entrenadores para el Mundial de Debate Escolar en Español (MED) de Chile 2025 y de Panamá 2026, donde la selección obtuvo su cuarto título consecutivo; el equipo que entrené llegó a cuartos de final y quedó entre los 10 mejores oradores del torneo.
  \item Coentrenador general de la Sociedad de Debate PUCP (desde julio de 2026) y entrenador del Colegio Alpamayo (desde junio de 2026).
  \item Duplas universitarias de República Dominicana, a distancia (desde 2024): en 2025, la primera dupla íntegramente dominicana en llegar a octavos de final del Campeonato Mundial Universitario de Debate en Español (CMUDE); en 2026, cuartos y octavos de final (CMUDE Ecuador).
  \end{puntos}}
\entrada[2025 -- 2026]{\cargo{Adjudicación}{torneos escolares y universitarios en Costa Rica, Ecuador, Perú y Venezuela}
  \detalle{Juez de la Gran Final del CMUDE 2026 (Ecuador) y jefe de adjudicación del Torneo Nacional de Debate Escolar (TNDE) 2025.}}
\entrada[2022 -- 2025]{\cargo{Como debatiente}{campeón del Campeonato Nacional Interuniversitario de Debate (CNDI 2025)}
  \detalle{Semifinalista del CMUDE Madrid 2022, cuartofinalista del CMUDE Panamá 2023 y subcampeón del Torneo Rosarista (Bogotá, 2023).}}
\end{entradas}

\seccion{Publicaciones e informes}
\begin{entradas}
\entrada[2026]{Monard, Encinas, Villena-Tagle y \textbf{Valverde} (equipo técnico). \textit{Análisis del elector peruano: aportes desde las ciencias del comportamiento}. Lima: Jurado Nacional de Elecciones y ProDiálogo, 85~pp.}
\entrada[2025]{\textbf{Valverde, L.} \textit{Desastres y brechas de implementación: estudio del proceso de reconstrucción de la infraestructura educativa en Piura entre julio de 2020 y enero de 2023, tras el Fenómeno de El Niño Costero}. Tesis de licenciatura con mención sobresaliente, PUCP. \enlace{https://hdl.handle.net/20.500.12404/32370}{Repositorio PUCP}.}
\entrada[2023]{Cabrales, S., Quiroz, S. J., Mamede, T., Barbosa, I. y \textbf{Valverde, L.} (equipo editor). \textit{Emergencia climática y derechos humanos: escrito sobre la opinión consultiva a la Corte Interamericana de Derechos Humanos}. Alianza de Clínicas Jurídicas Ambientales de Latinoamérica y el Caribe.}
\entrada[2023]{Seis columnas en coautoría en el boletín del IDEHPUCP (con I. Jave, M. Cépeda, A. Domínguez, E. Solís y A. Benites) sobre la búsqueda de personas desaparecidas, el RENADE, la enseñanza del conflicto armado interno, el cambio climático en los tribunales, la prevención de desastres y el estancamiento de la democracia en el Perú.}
\end{entradas}

\seccion{Reconocimientos y becas}
\begin{entradas}
\entrada[2024]{Beca de movilidad DARI-PUCP para el intercambio en Ghent University.}
\entrada[2023]{Primer puesto del Concurso de Investigación de la Facultad de Ciencias Sociales de la PUCP (pregrado), con una investigación sobre la reconstrucción posdesastre en Piura.}
\entrada[2019 -- 2024]{Beca 18 (PRONABEC): beca nacional por mérito que cubrió todo el pregrado.}
\end{entradas}

\seccion{Medios y paneles}
\begin{entradas}
\entrada[2026]{\cargo{Entrevistas}{\textit{El Salto} (España, julio), Radio Exitosa (junio y julio) y Canal 04 Catacaos, sobre la presidencia electa, los discursos de fraude y la coyuntura rumbo a las Elecciones Generales 2026.}}
\entrada[2026]{\cargo{Paneles}{\textit{La República}, «Beca 18: ¿solución o problema?» (julio), y Colegio de Politólogos del Perú, análisis del primer debate presidencial (marzo).}}
\end{entradas}

\seccion{Cursos de especialización}
\begin{entradas}
\entrada[2024 -- 2025]{\cargo{Q-LAB, PUCP}{métodos computacionales}
  \detalle{Innovaciones recientes en regresión discontinua para la inferencia causal (2025); Transformers in Large Language Models for Social Sciences and Public Policy (2025); introducción a \textit{machine learning} para las ciencias sociales y la gestión pública (2024); econometría aplicada avanzada (2024).}}
\entrada[2023]{\cargo{IDEHPUCP}{derechos humanos}
  \detalle{Derecho internacional humanitario, fuerzas públicas y derechos humanos (capacitación avanzada); introducción a la gestión ambiental con enfoque intercultural.}}
\entrada[2021]{\cargo{Korea Foundation}{Las economías de Corea del Sur y América Latina}}
\end{entradas}

\seccion{Métodos, herramientas e idiomas}
\begin{entradas}
\entrada[Cuantitativos]{Estadística descriptiva e inferencial; econometría e inferencia causal; \textit{machine learning} aplicado (XGBoost, SHAP); análisis de conglomerados; análisis espacial.}
\entrada[Cualitativos]{Estudios de caso; \textit{process tracing}; entrevistas semiestructuradas; análisis documental, de contenido y de discurso.}
\entrada[Análisis político]{Seguimiento legislativo y regulatorio; análisis institucional y de coyuntura; mapeo de actores.}
\entrada[Software]{R, mi herramienta principal (tidyverse, fixest, tidymodels); Python (scikit-learn, spaCy) y Stata, a nivel intermedio; Git y GitHub; LaTeX; KoboToolbox.}
\entrada[Idiomas]{Español, lengua materna. Inglés: First Certificate in English (B2); cursos de maestría en inglés en Ghent University.}
\end{entradas}

\seccion{Referencias}
\ifdefined\conreferencias
  \input{privado/referencias-es}
\else
  Disponibles a solicitud.
\fi

\end{document}
